% =============================================================================
% INTRODUCTION (Đặt vấn đề) — 1-2 pages
% Unnumbered chapter. Page numbering (1, 2, 3…) starts from this section.
% =============================================================================
\chapter*{INTRODUCTION}
\phantomsection
\addcontentsline{toc}{chapter}{INTRODUCTION}

Cervical cancer is a defining metric of global health inequality. Despite being entirely preventable through prophylactic human papillomavirus (HPV) vaccination, it claimed 342,000 lives in 2020---nearly 90\% of which occurred in low- and middle-income countries (LMICs) \cite{sung2021}. This disparity is not biological; it is fundamentally structural, rooted in the inequitable distribution of primary prevention \cite{arbyn2020}.

In Vietnam, the national vaccination strategy remains heavily reliant on out-of-pocket payments in the private sector. Consequently, while the WHO mandates a 90\% coverage target, actual uptake among Vietnamese women remains critically low \cite{bruni2023,nguyen2021}. More alarmingly, access to the vaccine---and even basic awareness of its existence---is strictly governed by socioeconomic fault lines: wealth, education, geography, and ethnicity \cite{tran2022,pham2020}.

To dismantle these structural barriers, health policy must be guided by precise econometric quantification. The Erreygers decomposition framework provides a mathematically rigorous method to isolate and rank the exact determinants driving health disparities \cite{erreygers2009}. While this methodology is established globally, it has yet to be applied to national-level HPV vaccination data in Vietnam.

Leveraging the 2020--2021 Vietnam Multiple Indicator Cluster Survey (MICS6), this thesis focuses on 4,102 women aged 15--29 \cite{unicef2022} to explicitly quantify the architecture of inequality in cervical cancer prevention among young Vietnamese women---the cohort most relevant to current HPV vaccination policy. By moving beyond simple descriptive statistics to robust econometric decomposition, this study aims to:

% --- Research objectives ---
\textbf{Research Objectives}
\begin{enumerate}
	\item Describe the weighted prevalence of HPV vaccine awareness and HPV vaccination among women aged 15--29 in Vietnam, stratified by sociodemographic characteristics.
	\item Quantify the magnitude of socioeconomic inequality in HPV vaccine awareness and vaccination using the concentration index.
	\item Identify and quantify the contribution of socioeconomic determinants to the observed inequality through Erreygers decomposition analysis.
\end{enumerate}

The findings are expected to provide actionable evidence for policymakers and program planners aiming to design equitable cervical cancer prevention strategies in Vietnam and comparable LMIC settings.

\clearpage
